\apendice{Documentación técnica de programación}

\section{Introducción}

En este apartado se describe la estructura técnica del sistema desarrollado, así como los aspectos necesarios para su instalación, ejecución y mantenimiento. Se detallan los módulos que componen cada uno de los addons, su organización interna y las principales funcionalidades implementadas.

El sistema está formado por dos addons diferenciados: el addon de recopilación de datos, implementado como un script ejecutable dentro de Blender, y el addon de análisis de datos, distribuido como un paquete instalable. Ambos componentes trabajan de forma complementaria mediante el uso de archivos CSV como medio de intercambio de información.

\section{Estructura de directorios}

El proyecto se divide en dos partes principales correspondientes a los dos addons desarrollados.

\subsection{Addon de recopilación de datos}

Este addon se implementa como un único archivo Python:

\begin{itemize}
    \item \textbf{DeteccionDatosV4.py}: script principal encargado de la captura de eventos, detección de cambios en la escena y registro de datos en archivos CSV.
\end{itemize}

Este enfoque permite una ejecución directa desde el editor de scripting de Blender, facilitando su uso en entornos académicos.

\subsection{Addon de análisis de datos}

Este addon se distribuye como un paquete estructurado en varios módulos:

\begin{itemize}
    \item \textbf{\_\_init\_\_.py}: archivo de inicialización del addon e integración con Blender.
    \item \textbf{metrics\_ui.py}: gestión de la interfaz de usuario y panel de visualización de métricas.
    \item \textbf{graphs.py}: generación de gráficos y representaciones visuales.
    \item \textbf{utils.py}: funciones auxiliares para procesamiento de datos.
    \item \textbf{operator.py}: definición de operadores de Blender para la interacción con el sistema.
    \item \textbf{site-packages/}: dependencias externas necesarias para el funcionamiento del addon.
\end{itemize}

Esta estructura modular facilita la mantenibilidad, escalabilidad y reutilización del código.

\section{Manual del programador}

El sistema ha sido desarrollado siguiendo un enfoque modular y orientado a eventos.

\subsection{Addon de recopilación}

El addon de recopilación se basa en el uso de \textit{handlers} de Blender para detectar eventos en la escena. Su funcionamiento principal consiste en:

\begin{itemize}
    \item Detectar cambios en la escena mediante eventos de Blender.
    \item Identificar el tipo de acción realizada (creación, modificación, eliminación).
    \item Capturar propiedades relevantes del objeto afectado.
    \item Registrar la información en archivos CSV.
\end{itemize}

Se han implementado mecanismos para evitar el registro de cambios redundantes, de forma que únicamente se almacenan modificaciones reales en el estado de los objetos.

Además, el sistema incorpora control de ejecución mediante funcionalidades de inicio y parada (Start/Stop), permitiendo al usuario gestionar la captura de datos.

\subsection{Addon de análisis}

El addon de análisis está diseñado de forma modular, separando claramente la lógica de procesamiento de datos de la interfaz de usuario.

Sus componentes principales son:

\begin{itemize}
    \item \textbf{Carga de datos}: lectura de archivos CSV generados por el addon de recopilación.
    \item \textbf{Procesamiento de métricas}: cálculo de estadísticas y métricas a partir de los datos registrados.
    \item \textbf{Visualización}: representación de resultados mediante gráficos y paneles dentro de Blender.
\end{itemize}

El sistema permite analizar la información de forma flexible, ofreciendo diferentes tipos de visualización y facilitando la interpretación de los datos.

\section{Compilación, instalación y ejecución del proyecto}

\subsection{Addon de recopilación}

El addon de recopilación se ejecuta directamente desde Blender siguiendo estos pasos:

\begin{enumerate}
    \item Abrir Blender.
    \item Acceder al espacio de trabajo \textit{Scripting}.
    \item Cargar el archivo \texttt{DeteccionDatosV4.py}.
    \item Activar la opción \textit{Register}.
    \item Ejecutar el script.
    \item Guardar el archivo \texttt{.blend}.
\end{enumerate}

Una vez ejecutado, el sistema solicita el consentimiento del usuario para la recopilación de datos. A partir de ese momento, la captura de eventos se realiza automáticamente durante el uso del proyecto.

\subsection{Addon de análisis}

El addon de análisis se instala como un addon estándar de Blender:

\begin{enumerate}
    \item Comprimir el addon en formato \texttt{.zip}.
    \item Acceder a \textit{Edit > Preferences > Add-ons}.
    \item Seleccionar \textit{Install} e importar el archivo \texttt{.zip}.
    \item Activar el addon.
\end{enumerate}

Una vez instalado, el usuario puede acceder al panel de análisis y cargar los datos generados para su procesamiento.

\section{Pruebas del sistema}

Se han realizado diferentes pruebas para validar el correcto funcionamiento del sistema.

\subsection{Pruebas del addon de recopilación}

\begin{itemize}
    \item Registro correcto de creación de objetos.
    \item Detección de modificaciones en \textit{Edit Mode}.
    \item Captura de operaciones como duplicado (Shift+D), pegado (Ctrl+V) y eliminación.
    \item Verificación del funcionamiento de los controles Start/Stop.
    \item Comprobación de que no se registran cambios redundantes.
    \item Validación del almacenamiento correcto en archivos CSV.
\end{itemize}

\subsection{Pruebas del addon de análisis}

\begin{itemize}
    \item Carga correcta de archivos CSV.
    \item Verificación del cálculo de métricas.
    \item Validación de la representación gráfica (escalas, ejes y datos).
    \item Comprobación de la coherencia de la información mostrada en pantalla.
    \item Pruebas de funcionamiento de distintos tipos de gráficos.
\end{itemize}

Estas pruebas permiten garantizar que el sistema funciona de forma consistente, tanto en la captura de datos como en su posterior análisis.

El sistema ha sido validado de forma incremental durante su desarrollo, asegurando la coherencia entre los datos capturados y los resultados obtenidos en el análisis.